\chapter{\IfLanguageName{dutch}{Gebruikerstesting en evaluatie}{Usertesting and evaluation}}%
\label{ch:gebruikerstesting}

\section{Opzet van de testing}
\label{sec:opzet}

De gehoste applicatie werd verspreid onder huisartsen en praktijkmedewerkers via een uitnodigingsmail. 
32 praktijken binnen Eerstelijnszone Aalst en omligende regio's werden hiervoor gecontacteerd.
De deelnemers werden gevraagd om de interface zelfstandig te doorlopen en na afloop een vragenlijst in te vullen bestaande uit twee onderdelen:

\begin{itemize}
    \item De gestandaardiseerde System Usability Scale (SUS).
    \item Een open feedback gedeelte.
\end{itemize}

Voor de afname van deze evaluatievragenlijst werd gebruikgemaakt van \textbf{Google Forms}. 
Deze keuze werd gemaakt om de antwoorden centraal en gestructureerd te verzamelen en de resultaten eenvoudig te kunnen exporteren naar een spreadsheet voor verdere verwerking en analyse. 
De spreadsheet werd gebruikt voor het berekenen van de individuele en gemiddelde SUS-scores, evenals de mediaan van de SUS-scores. 

Naast de vragenlijst werd gebruikgemaakt van \textbf{Microsoft Clarity} voor het verzamelen van gebruikersinteracties en tijdsmetingen. 
Via sessie-opnames en interactiedata konden onder meer de registratietijd en de totale invultijd van de vragenlijst worden geanalyseerd. 
Daarnaast maakte Microsoft Clarity het mogelijk om gebruikersgedrag te observeren via klik- en scrollgedrag, wat aanvullende inzichten bood in mogelijke gebruiksproblemen binnen de interface.


De System Usability Scale bestaat uit tien stellingen die worden beoordeeld op een vijfpunts Likertschaal en levert een score op tussen 0 en 100.

\begin{enumerate}
    \item Ik denk dat ik deze applicatie graag regelmatig zou gebruiken.
    \item Ik vond de applicatie onnodig complex.
    \item Ik vond de applicatie gemakkelijk te gebruiken.
    \item Ik denk dat ik de ondersteuning van een technisch persoon nodig zou hebben om deze applicatie te kunnen gebruiken.
    \item Ik vond de verschillende functies in deze applicatie goed geïntegreerd.
    \item Ik vond dat er veel inconsistenties in de applicatie zaten.
    \item Ik kan me voorstellen dat de meeste mensen deze applicatie snel onder de knie zouden krijgen.
    \item Ik vond de applicatie onnodig ingewikkeld om te gebruiken.
    \item Ik voelde me erg zeker bij het gebruiken van deze applicatie.
    \item Ik moest veel leren voordat ik met deze applicatie aan de slag kon.
\end{enumerate}

Het open feedbackgedeelte liet deelnemers vrij aangeven wat goed werkte, wat voor verbetering vatbaar was en welke elementen zij als verwarrend of onduidelijk ervaarden.
Dit gedeelte bestond uit 5 vragen:

\begin{enumerate}
    \item Wat vond je het meest waardevol aan de applicatie?
    \item Wat vond je het minst prettig of verwarrend?
    \item Mis je bepaalde functies die je zou verwachten?
    \item Zou je deze applicatie gebruiken in je dagelijkse praktijk? Waarom wel/niet?
    \item Heb je nog andere opmerkingen of suggesties?
\end{enumerate}

De combinatie van subjectieve evaluatie via de SUS-score en objectieve interactiedata via Microsoft Clarity maakte het mogelijk om de gebruiksvriendelijkheid van de interface vanuit meerdere perspectieven te beoordelen.

\section{Deelnemers}
\label{sec:deelnemers}

In totaal namen 12 huisartsen en/of praktijkmedewerkers deel aan de testing. 
Dit komt neer op een responsegraad van 35,3\%.

De deelnemers werden gerekruteerd via e-mailcommunicatie naar huisartsenpraktijken binnen Eerstelijnszone Aalst en omliggende regio's. 
Tijdens de rekrutering bleek dat persoonlijke contactkanalen een duidelijke invloed hadden op de responsgraad. 
Praktijken waarbij rechtstreeks contact mogelijk was met een arts of medewerker reageerden merkbaar vaker dan praktijken die uitsluitend via een algemeen praktijkadres werden gecontacteerd. 
Daarnaast leidde ook communicatie via tussenpersonen of bestaande professionele contacten tot een hogere respons.

De meerderheid van de ontvangen antwoorden was afkomstig van groepspraktijken en in beperktere mate van duopraktijken. 
Er werden geen bruikbare responses ontvangen vanuit solopraktijken, waardoor deze categorie niet vertegenwoordigd is binnen de steekproef. 
Dit vormt een beperking van het onderzoek, aangezien de gebruikservaring en organisatorische context van solopraktijken mogelijk verschillen van grotere praktijkstructuren.

De steekproef kan bijgevolg niet als volledig representatief worden beschouwd voor alle types huisartsenpraktijken binnen de eerstelijnszorg. 
Desondanks biedt de testing waardevolle inzichten in de bruikbaarheid van de ontwikkelde interface binnen een reële zorgcontext.

\section{SUS-resultaten}
\label{sec:susresultaten}

De System Usability Scale (SUS) werd gebruikt om de ervaren gebruiksvriendelijkheid van de ontwikkelde interface te evalueren. 
De SUS-score wordt berekend volgens de gestandaardiseerde methode waarbij de antwoorden op de tien stellingen worden omgerekend naar een score op 100.

De gemiddelde SUS-score over alle deelnemers bedraagt 73,96, met een mediaanscore van 73,75. 
De laagste geregistreerde score bedroeg 50, terwijl de hoogste score 90 bedroeg.

\begin{table}[H]
    \centering
    \begin{tabular}{|l|c|}
        \hline
        \textbf{Metriek} & \textbf{Score} \\
        \hline
        Gemiddelde SUS-score & 73,96 \\
        Mediaan SUS-score & 73,75 \\
        Minimumscore & 50 \\
        Maximumscore & 90 \\
        \hline
    \end{tabular}
    \caption{Overzicht van de SUS-resultaten}
    \label{tab:susresultaten}
\end{table}

Volgens de benchmarkwaarden uit de literatuur geldt een SUS-score van 68 als het gemiddelde voor commerciële systemen, terwijl scores boven 75 worden beschouwd als indicatief voor een goede usability en scores boven 85 als uitstekend. 
De bekomen resultaten wijzen bijgevolg op een bovengemiddelde gebruiksvriendelijkheid van de ontwikkelde interface.

De beperkte spreiding tussen de gemiddelde en mediaanscore suggereert bovendien dat de resultaten relatief consistent verdeeld zijn over de verschillende deelnemers. 
Een verdere analyse van de individuele scores en kwalitatieve feedback wordt besproken in het volgende hoofdstuk.

\section{Tijdsmetingen}
\label{sec:tijdsmetingen}

Naast de SUS-score werden ook tijdsmetingen geanalyseerd op basis van Microsoft Clarity sessie-opnames. 
Hierbij werden zowel de registratietijd als de totale invultijd van de evaluatievragenlijst gemeten.

De resultaten tonen een gemiddelde, mediaan, minimum en maximum voor beide metrieken. 
Deze waarden geven inzicht in de efficiëntie waarmee deelnemers de applicatie konden doorlopen en de vragenlijst konden invullen. 
De exacte waarden zijn opgenomen in onderstaande tabel.

\begin{table}[H]
    \centering
        \begin{tabular}{|l|c|c|}
            \hline
            \textbf{Metriek} & \textbf{Registratietijd} & \textbf{Invultijd} \\
            \hline
            Gemiddelde & 4:16 & 3:17 \\
            Mediaan & 3:15 & 2:28 \\
            Minimum & 1:16 & 1:21 \\
            Maximum & 11:34 & 9:28 \\
            \hline
        \end{tabular}
    \caption{Overzicht van de tijdsmetingen gebaseerd op Microsoft Clarity}
    \label{tab:tijdsmetingen}
\end{table}

\section{Kwalitatieve feedback}
\label{sec:kwalitatievefeedback}

Naast de kwantitatieve SUS-score werd kwalitatieve feedback verzameld via het open gedeelte van de vragenlijst. 
Deze feedback werd thematisch gecodeerd op basis van terugkerende patronen en gegroepeerd in overkoepelende thema's.

De analyse resulteerde in vier dominante thema's: (1) inhoudelijke meerwaarde van de vragenlijst, (2) gebruiksvriendelijkheid en interfaceproblemen, (3) functionele uitbreidingen en automatisatie, en (4) adoptie en contextuele relevantie binnen de eerstelijnszorg.

\subsection{Inhoudelijke meerwaarde}
\label{subsec:inhoudelijkeMeerwaarde}

Het eerste terugkerende thema betreft de inhoudelijke meerwaarde van de vragenlijst. 
Verschillende deelnemers gaven aan dat de bevraagde informatie relevant en niet standaard beschikbaar is binnen bestaande systemen. 
Vooral de bevraging van praktijkkenmerken en organisatiegegevens werd als nuttig beschouwd. 
Tegelijk werd opgemerkt dat bepaalde vragen potentieel uitgebreid kunnen worden, bijvoorbeeld rond werkbelasting en praktijkorganisatie.

\subsection{Gebruiksvriendelijkheid en interfaceproblemen}
\label{subsec:gebruiksvriendelijkheidinterfaceproblemen}

Een tweede thema betreft de gebruikservaring van de interface. 
Hoewel de applicatie over het algemeen als vlot invulbaar werd ervaren, werd feedback gegeven op visuele aspecten zoals de hoeveelheid witruimte en de duidelijkheid van bepaalde UI-elementen, zoals de ja/nee-schuivers. 
Dit wijst op verbeterpunten in visuele hiërarchie en componentkeuze binnen de interface.


\subsection{Functionaliteit en automatisatie}
\label{subsec:functionaliteitAutomatisatie}

Een derde thema betreft gewenste uitbreidingen op vlak van functionaliteit. 
Deelnemers suggereerden onder meer automatische invulling van praktijkgegevens, bijvoorbeeld via koppeling met bestaande registers zoals RIZIV-data. 
Daarnaast werd het ontbreken van automatische opslag bij navigatie tussen verschillende delen van de vragenlijst als een beperking ervaren.

\subsection{Adoptie en context}
\label{subsec:adoptiecontext}

Een vierde thema betreft de contextuele relevantie en adoptie van de tool. 
Verschillende deelnemers gaven aan dat de meerwaarde van de applicatie afhankelijk is van de bestaande organisatie binnen de praktijk en van reeds gebruikte alternatieven. 
In sommige gevallen werd verwezen naar bestaande oplossingen zoals regionale huisartsenzoekers, waardoor de nood aan een bijkomende tool beperkter werd ingeschat.

\subsection{Overzicht}
\label{subsec:overzicht}

Overkoepelend tonen de resultaten aan dat de applicatie vooral wordt gewaardeerd om haar inhoudelijke relevantie en vlotte invulbaarheid, maar dat er duidelijke verbeterpunten zijn op vlak van gebruikersinterface en automatisatie. 
Daarnaast blijkt dat de adoptie sterk contextafhankelijk is en beïnvloed wordt door bestaande digitale oplossingen binnen de eerstelijnszorg.

\section{Beperkingen tijdens de testing}
\label{sec:beperkingen}

Tijdens de uitvoering van de gebruikerstesting werden enkele technische en methodologische beperkingen vastgesteld die een impact hebben op de volledigheid van de verzamelde data.

Ten eerste werd een compatibiliteitsprobleem vastgesteld bij het gebruik van de applicatie op Firefox en bepaalde Apple-apparaten. 
Hierbij werden de authenticatiegegevens niet correct opgeslagen, waardoor foutieve of mislukte requests naar de backend ontstonden. 
Het compatibiliteitsprobleem dat zich voordeed bij het gebruik van Firefox en bepaalde Apple-apparaten kon tijdens de testfase worden geïdentificeerd en opgelost. 
De oorzaak bleek te liggen in de configuratie van de authenticatieflow, waarbij de nodige OAuth-scopes niet volledig correct waren ingesteld.

Door het toevoegen van de ontbrekende parameters binnen de Auth0-configuratie in de Angular applicatie werd het probleem opgelost:

\begin{minted}{typescript}
    audience: environment.auth0.authorizationParams.audience,
    scope: 'openid profile email offline_access',
\end{minted}

Hierdoor werden de authenticatiegegevens correct opgeslagen en konden requests naar de backend opnieuw succesvol worden uitgevoerd.

Ten tweede trad er een configuratieprobleem op bij het herstarten van de API-server, waarbij de omgevingsvariabelen (secrets) niet correct werden geïnitialiseerd. 
Hierdoor bleef de API tijdelijk offline totdat de server manueel werd herstart.

Ten slotte bleken er voor twee deelnemers geen Microsoft Clarity-opnames beschikbaar te zijn, waardoor voor deze cases geen tijdsmetingen konden worden geregistreerd. 
Dit resulteert in een beperkte dataset voor de analyse van registratietijd en invultijd.